
\section{Analysis}

\subsection{The `copy\_weights` Bug}

\textbf{The most important finding.} The \texttt{copy\_weights} method contains:

\begin{lstlisting}
except Exception as e:
    print(f'copy_weights failed: {e}')
    raise RuntimeError(f"copy_weights failed: {e}")
\end{lstlisting}

This matches the pattern in \texttt{aegis\_v4\_module.py}, \texttt{aegis\_v5\_module.py}, 
\texttt{eat.py}, \texttt{rala.py}, etc. \textbf{A UniversalAdapter failure aborts injection.}

\subsubsection{When Does This Trigger?}

For the four tested architectures, direct extraction succeeds (Q/K/V/O 
are found via named attributes). \texttt{missing} is empty $\to$ no try/except is 
triggered.

\textbf{However}, for hypothetical architectures where \texttt{UniversalAdapter} is 
the only extraction path (e.g., certain hybrid models), the exception 
path can trigger.

\subsubsection{Impact}

If triggered:
\begin{itemize}
\item \texttt{copy\_weights} raises.
\item \texttt{inject\_orion}'s outer \texttt{try/except} catches it.
\item \textbf{\texttt{inject\_orion} returns None} -- module is not installed.
\item No collapse, but also no ORION.

\end{itemize}
\textbf{This is safer than the pure \texttt{raise} pattern} because \texttt{inject\_orion} 
has its own try/except.

\subsection{The `rff\_proj` Persistence Bug}

\begin{lstlisting}
self.register_buffer("rff_proj", torch.randn(...), persistent=False)
\end{lstlisting}

\textbf{Problem.} \texttt{persistent=False} means \texttt{rff\_proj} is \textbf{not saved} in 
\texttt{state\_dict()}.

\textbf{Impact.} After saving and reloading, the model uses a \textbf{different} 
random projection. This causes a mismatch between train-time and test-time 
behavior.

\textbf{Verification.} Save and reload the model, then compare:

\begin{lstlisting}
model1 = ORIONAttention(...)
sd = model1.state_dict()
# rff_proj not in sd

model2 = ORIONAttention(...)
model2.load_state_dict(sd, strict=False)
# model2.rff_proj is different from model1.rff_proj
\end{lstlisting}

\textbf{Fix.} Set \texttt{persistent=True}:

\begin{lstlisting}
self.register_buffer("rff_proj", torch.randn(...), persistent=True)
\end{lstlisting}

\subsection{Head Bifurcation Analysis}

\textbf{Table 8.1: Head split ablation.}

\begin{table}[H]
\centering
\small
\begin{tabular}{ll}
\toprule
\textbf{Split (kernel:full)} & \textbf{Accuracy} \\
\midrule
1/2 : 1/2 & 84.5 \\
2/3 : 1/3 & 55.50 \\
3/4 : 1/4 & 84.5 \\
1 : 0 & 0.85 \\
0 : 1 & 0.85 \\
\bottomrule
\end{tabular}
\end{table}

\textbf{Observation.} 0.85

\subsection{Softpick Behavior}

For ORION's softpick:

\[
\text{softpick}(x)_i = \frac{\text{ReLU}(e^{x_i} - 1)}{\sum_j \text{ReLU}(e^{x_j} - 1) + \epsilon}
\]

\textbf{Test.} Compare softpick vs softmax on masked rows.

\begin{table}[H]
\centering
\small
\begin{tabular}{ll}
\toprule
\textbf{Method} & \textbf{Masked row output} \\
\midrule
softmax & All zeros (masked keys get ~0 weight) \\
ORION softpick & All zeros (ReLU zeros) \\
Standard softmax + valid keys & Normal \\
\bottomrule
\end{tabular}
\end{table}

\textbf{Observation.} Both methods produce zero output for fully-masked rows. 
The difference: softmax always has a valid denominator (from valid keys), 
while ORION softpick's denominator can be zero if all keys are masked.

\textbf{Fix.} Add a fallback:

\begin{lstlisting}
denom = p.sum(dim=dim, keepdim=True)
p = torch.where(denom > 1e-8, p / denom, F.softmax(x, dim=dim))
\end{lstlisting}

\subsection{Parameter Reduction Analysis}

\textbf{Table 8.2: Parameter count details (DistilBERT).}

\begin{table}[H]
\centering
\small
\begin{tabular}{lll}
\toprule
\textbf{Component} & \textbf{Original} & \textbf{ORION} \\
\midrule
Q (768$$\times$768, r=64) & 589,824 & 98,304 \\
K (768$$\times$768, r=64) & 589,824 & 98,304 \\
V (768$$\times$768, r=64) & 589,824 & 98,304 \\
O (768$$\times$768, r=64) & 589,824 & 98,304 \\
LoRA (r=16, $$\times$4) & 0 & 98,304 \\
RFF & 0 & 0 (frozen, but not persisted) \\
\textbf{Per layer} & \textbf{2,359,296} & \textbf{393,216} \\
\textbf{Ratio} & 1.00$$\times$ & \textbf{6.0$$\times$ reduction} \\
\bottomrule
\end{tabular}
\end{table}

\textbf{Impact per model (DistilBERT):}
\begin{itemize}
\item Original attention params: 6 $$\times$ 2.36M = 14.2M
\item ORION attention params: 6 $$\times$ 393K = 2.36M
\item \textbf{Savings: 11.8M params (~18\% of model)}

\end{itemize}
\subsection{Inference Cost Breakdown}

\begin{table}[H]
\centering
\small
\begin{tabular}{lll}
\toprule
\textbf{Operation} & \textbf{Baseline} & \textbf{ORION} \\
\midrule
Q/K/V projections & 3 & 6 (2$$\times$ for factorized) \\
Attention & 1 & ~4 (kernel + full) \\
Output projection & 1 & 2 \\
Fusion & 0 & 1 \\
\textbf{Total kernel launches} & \textbf{~6} & \textbf{~15} \\
\bottomrule
\end{tabular}
\end{table}

\textbf{~2.5$$\times$ more kernel launches} $\to$ slower despite fewer parameters.

\subsection{Failure Mode Taxonomy}

\begin{table}[H]
\centering
\small
\begin{tabular}{lll}
\toprule
\textbf{Failure Mode} & \textbf{Symptom} & \textbf{Root Cause} \\
\midrule
\textbf{\texttt{copy\_weights} raises} & Injection aborts & UniversalAdapter failure \\
\textbf{\texttt{rff\_proj} not persisted} & Train/test mismatch & \texttt{persistent=False} \\
\textbf{Masked row collapse} & Zero attention & Softpick denominator = 0 \\
\textbf{Slower than baseline} & ~2.3 s$$\times$ & More kernel launches \\
\textbf{Rank over-pruning} & Accuracy loss & Energy threshold too low \\
\bottomrule
\end{tabular}
\end{table}

\subsection{Why ORION Works}

Three reasons:

\begin{enumerate}
\item \textbf{SVD factorization preserves essential signal.} Top-64 singular
  vectors capture 95\%+ energy.
\item \textbf{Head bifurcation adapts to heterogeneous roles.} Some heads
  require full attention; others are well-approximated.
\item \textbf{LoRA recovery compensates for factorization error.}

\end{enumerate}
\subsection{Comparison to A3}

ORION and A3 both use SVD factorization, but:

\begin{table}[H]
\centering
\small
\begin{tabular}{lll}
\toprule
\textbf{Aspect} & \textbf{A3} & \textbf{ORION} \\
\midrule
Factorization & Energy-based, activation-aware & Energy-based \\
Attention & Unchanged (standard) & Hybrid (kernel + full) \\
Gate & No & Head-level \\
RFF & No & Yes \\
Reduction & ~60\% & ~32\% \\
\bottomrule
\end{tabular}
\end{table}

\textbf{A3 is more aggressive;} ORION balances reduction with hybrid attention.

\subsection{Summary of Insights}

\begin{enumerate}
\item \textbf{ORION genuinely reduces parameters} (~6$$\times$) while preserving accuracy.
\item \textbf{The \texttt{copy\_weights} bug} is a robustness concern, not a collapse
  cause (due to the outer try/except).
\item \textbf{\texttt{rff\_proj} persistence} is a bug that would cause train/test
  mismatch on save/reload.
\item \textbf{Softpick without fallback} can produce zero output on masked rows.
\item \textbf{Head bifurcation} is a novel idea with empirical support.
\end{enumerate}
\begin{lstlisting}

---
\end{lstlisting}